% !TEX program = xelatex
\documentclass{article}
\usepackage{ctex}
\usepackage[margin=1in]{geometry}
\usepackage{terminalcode}
% Set terminal theme: 'dark' (default) or 'light'
\terminalcodetheme{dark}

% ============================================================================
% USAGE GUIDE (auto-generated by cmdlog2tex)
%
% 1. Inline terminal code block:
%      \begin{termcode}[<language>]{<title>}
%        your code or log content here
%      \end{termcode}
%
%    - <language>: optional (e.g., bash, python, text). Default: text.
%    - <title>: required (appears in box header).
%
%    Example:
%      \begin{termcode}[bash]{Shell Session}
%        $ ls -la
%        $ echo "Hello"
%      \end{termcode}
%
% 2. Include external file with same styling:
%      \terminput[<language>]{<title>}{<file-path>}
%
%    This is the safe, recommended way to include files (avoids nesting issues).
%    Alias \termcodefileinput is also available.
%
%    Example:
%      \terminput[bash]{Build Script}{build.sh}
%
% 3. Embed colored text inside listings (using escapeinside=«»):
%      «\ansicolor{1}{32}some green bold text»
%
%    - First argument: ANSI text attribute (e.g., "1" for bold, "" for normal).
%    - Second argument: ANSI color code (30–37 or 90–97).
%    - \ac is a short alias for \ansicolor.
%
%    Valid color codes:
%      30–37: black, red, green, yellow, blue, magenta, cyan, white
%      90–97: bright black/gray, bright red, ..., bright white
%
%    Note: \verb cannot be styled directly. Use escape markers «...» instead.
%
% 4. Theme switching:
%      \tctheme{dark}   % or \terminalcodetheme{light}
%
% 5. To typeset literal « or » in output:
%      Copy-paste the characters directly, or ensure your editor supports UTF-8.
%      They are Unicode U+00AB and U+00BB.
%
% ============================================================================
\title{CMDLOG2TEX Demo}
\author{LoveElysia1314}
\date{November 1, 2025}

\begin{document}

\maketitle

\section{Dark Theme with ANSI Colors}

Default dark theme showing terminal output with ANSI color preservation:

\terminput[bash]{Bash Session (Dark)}{test_bash_demo.color.txt}

\newpage

\section{Light Theme with ANSI Colors}

Same output rendered with light theme:

\terminalcodetheme{light}

\terminput[bash]{Bash Session (Light)}{test_bash_demo.color.txt}

\newpage

\section{Plain Text Output}

Plain text version without ANSI color codes:

\terminalcodetheme{dark}

\terminput[text]{Bash Session (Plain)}{test_bash_demo.plain.txt}

\newpage

\section{Theme Switching}

To switch themes in your document, use:
\begin{verbatim}
\terminalcodetheme{dark}    % Dark theme (default)
\terminalcodetheme{light}   % Light theme
\end{verbatim}

\section{Including Files}

To include terminal output files in your document:
\begin{verbatim}
\terminput[<language>]{<title>}{<file-path>}
\end{verbatim}

Parameters:
\begin{itemize}
  \item \texttt{<language>}: Optional language tag (bash, python, text, etc.)
  \item \texttt{<title>}: Display title for the output box
  \item \texttt{<file-path>}: Relative or absolute path to the file
\end{itemize}

\end{document}
